%%
%% commutative-algebra.tex
%% 
%% Made by Alex Nelson <pqnelson@gmail.com>
%% Login   <alex@lisp>
%% 
%% Started on  2025-11-11T16:19:29-0800
%% Last update 2025-11-11T16:19:29-0800
%% 

\chapter{Commutative Algebra}

\begin{remark}[References]
There are quite a few references for commutative algebra.
Atiyah and Macdonald~\cite{atiyah1969introduction} remains a classic,
Eisenbud~\cite{eisenbud1994commutative} contains a number of good
examples and exercises. Matsumara~\cite{matsumura1970commutative}
is also quite good.
\end{remark}

\begin{definition}[Spectrum of a ring]\mml{topzari:def 5}
Let $R$ be a commutative ring. We define the \define{Spectrum} of $R$
to be the subset-family $\Spec(R)$ equal to $\{I\ideal R\mid I\mbox{ is prime}\}$.
% Stacks: 00DY
\end{definition}

\begin{definition}[Maximal spectrum of a ring]\mml{topzari:def 7}
Let $R$ be a commutative ring.
We define the \define{Maximal Spectrum} of $R$ to be the set $\MSpec(R)$ of all
maximal ideals of $R$.
\end{definition}

\begin{theorem}
If $R$ is a commutative ring and $I\ideal R$ such that $I\neq R$,
then $R$ has a maximal ideal $\mathfrak{m}$ such that $I\subset\mathfrak{m}$.
\end{theorem}

\begin{corollary}
Every non-unit of a commutative ring lies in a maximal ideal.
\end{corollary}

\begin{proof}
If $x$ is a non-unit, then $(x)\neq R$. Then apply the previous theorem.
\end{proof}

\begin{proposition}
Let $R$ be a ring with a maximal ideal $\mathfrak{m}$. If every
element of $1+\mathfrak{m}$ is a unit, then $R$ is a local ring.
\end{proposition}

\begin{proof}
Let $x\in R\setminus\mathfrak{m}$.
Since $\mathfrak{m}$ is maximal, the smallest ideal containing both
$\mathfrak{m}$ and $x$ is $R$. Then $1=ax+y$ for some $a\in R$ and
$y\in\mathfrak{m}$. Then $ax=1-y$ is a unit by assumption, so
$\mathfrak{m}$ contains all the non-units.
\end{proof}

\begin{proposition}[Ideal is prime iff quotient ring is integral domain]
Let $R$ be a unital commutative ring. Let $I\ideal R$.
Then $I$ is a prime ideal of $R$ if and only if the quotient ring
$R/I$ is an integral domain.
\end{proposition}

\begin{proof}
\forwardproof\ Let $I$ be a prime ideal. We want to show: for any $(x+I)\in R/I$ and
$(y+I)\in R/I$ such that $(x+I)(y+I)=xy+I=0_{R/I}$, then we have $x+I=0_{R/I}$
or $y+I=0_{R/I}$.

The zero of $R/I$ is $0_{R/I}=I$. Therefore $xy+I=0_{R/I}$ implies
$xy\in I$. Since $I$ is a prime ideal, either $x\in I$ or $y\in I$.
Without loss of generality, we may assume $x\in I$. Then $x+I=I=0_{R/I}$.
Hence the claim.

\backwardproof\ Assume $R/I$ is an integral domain.
Let $x,y\in R$ be such that $xy\in I$. Then $0_{R/I}=I=xy+I=(x+I)(y+I)$.
Since $R/I$ is an integral domain, it follows that either
$x+I=0_{R/I}$ or $y+I=0_{R/I}$. If $x+I=0_{R/I}$, then $x\in I$.
If $y+I=0_{R/I}$, then $y\in I$. Hence either $x\in I$ or $y\in I$,
and so by definition $I$ is a prime ideal.
\end{proof}

\begin{proposition}[Maximal ideals of commutative rings are prime]
Let $R$ be a commutative ring.
If $\mathfrak{m}\ideal R$ is a maximal ideal, then $\mathfrak{m}$ is a
prime ideal of $R$.
\end{proposition}

\begin{proof}
The quotient $R/\mathfrak{m}$ is a field (and in particular, an
integral domain). But an ideal $I$ is prime iff $R/I$ is an integral
domain. Hence $\mathfrak{m}$ must be prime.
\end{proof}

\begin{definition}[Local ring]\mml{topzari:def 9}
Let $R$ be a commutative ring where $1\neq0$.
We say $R$ is \define{Local} if there is exactly one maximal ideal of $R$.
(This could be defined for any nonempty double loop structure $R$.)
\end{definition}

\begin{definition}[Semi-local rings]\mml{topzari:def 10}
Let $R$ be a commutative ring where $1\neq0$.
We say $R$ is \define{Semi-Local} if there are only finitely many
maximal ideals of $R$. Equivalently, $\MSpec(R)$ is a finite set.
\end{definition}

\begin{definition}[Nilradical of a ring]\mml{topzar1:def 13}
Let $A$ be a nondegenerate commutative ring.
We define the \define{Nilradical} of $A$ to be the ideal
$\NilRadical(A):=\{a\in A\mid a\mbox{ is nilpotent}\}$. 
\end{definition}

\begin{theorem}\mml{topzar1:18}
Let $A$ be a nondegenerate commutative ring.
Let $I\ideal A$ be an ideal.
Then $\Radical{I}=\Intersects\{P\in\Spec(A)\mid I\subset P\}$.
\end{theorem}

\begin{theorem}\mml{topzar1:19}
Let $A$ be a nondegenerate commutative ring.
Then $\NilRadical(A)=\Intersects\Spec(A)$.
\end{theorem}

\begin{definition}[Reduced rings have no nilradical]\mml{field_16:20}
Let $R$ be a nondegenerate commutative ring.
We say $R$ is \define{Reduced} when $\NilRadical{R}=\{0\}$.
\end{definition}

\begin{theorem}
Let $A$ be an arbitrary commutative ring. Then the quotient ring
\begin{equation*}
R = A/\bigcap\Spec(A)
\end{equation*}
is a reduced ring.
\end{theorem}

\begin{definition}[Jacobson radical of a ring]
Let $R\neq0$ be an arbitrary ring (not necessarily commutative).
Then the \define{Jacobson Radical} of $R$ is the ideal
\begin{equation}
\JacobsonRadical(R):=\{r\in R\mid rM=0\mbox{ for all }M\mbox{ simple}\}
\end{equation}
consisting of all elements $r\in R$ which annihilate all simple
modules over $R$.
\end{definition}

\begin{proposition}[Jacobson radical of a commutative ring]
Let $R\neq0$ be a commutative ring. Then the Jacobson Radical of $R$
is the ideal $\JacobsonRadical(R)=\bigcap\MSpec(R)$ obtained by taking
the intersection of all maximal ideals of $R$.
\end{proposition}

The proposition holds for noncommutative rings, if we work with
maximal left ideals.

\begin{proof}
\textsc{Claim 1:} $\JacobsonRadical(R)\subset\bigcap\MSpec(R)$.
Let $x\in\JacobsonRadical(R)$ and $\mathfrak{m}$ be a maximal ideal.
Then $R/\mathfrak{m}$ is a simple module, so $x(R/\mathfrak{m})=0$.
Hence $x\in\mathfrak{m}$.

\textsc{Claim 2:} $\JacobsonRadical(R)\supset\bigcap\MSpec(R)$.
Let $x$ be in every maximal ideal, let $M$ be a simple module over $R$.
Then for every nonzero $v\in M$, $Rv=M$ and $\Annihilator(v)$ is a
maximal ideal. Then $xv=0$. Then $xM=0$, hence $x\in\JacobsonRadical(R)$.
\end{proof}

\begin{proposition}
Let $R$ be a commutative ring. If $x\in\JacobsonRadical(R)$,
then $1+x$ is a unit in $R$.
\end{proposition}

The previous proposition is stated in Matsumara~\cite{matsumura1970commutative}
after Definition 1L.

\begin{proof}
Let $x\in\JacobsonRadical(R)$. Assume for contradiction $1+x$ is not a
unit in $R$.
\end{proof}

\begin{proposition}
Let $I\ideal R$ be an ideal. If $1+x$ is a unit for all $x\in I$,
then $I\subset\JacobsonRadical(R)$.
\end{proposition}

The previous proposition is stated in Matsumara~\cite{matsumura1970commutative}
after Definition 1L.

\begin{proof}
Assume $1+x$ is a unit for all $x\in I$.
Assume for contradiction $\mathfrak{m}$ is a maximal ideal and
$I\nsubset\mathfrak{m}$.
Then $I+\mathfrak{m}=R$.
Then $1=x+y$ for some $x\in I$ and $y\in\mathfrak{m}$.
Then $y=1-x$ is not a unit. Hence contradiction.
\end{proof}

\section{Localization}

\begin{definition}[Multiplicative set]\mml{c0sp1:def 4}
Let $M$ be a nonempty multiplicative loop with zero.
Let $S$ be a subset of $M$.
We say $S$ is \define{Multiplicatively-Closed} (or just \emph{multiplicative})
if
\begin{enumerate}
\item $1\in S$, and
\item for all $u,v\in S$, we have $uv\in S$.
\end{enumerate}
\end{definition}

\begin{proposition}[Partial converse of prime ideal, to construct prime ideals]
Let $R$ be a commutative ring, let $U\subset R$ be a
multiplicatively-closed subset. If $I\properideal R$ is an ideal maximal
among those disjoint with $U$, then $I$ is prime.
\end{proposition}

This partial converse is due to
Eisenbud~\cite[Prop.~2.11]{eisenbud1994commutative}.

\begin{proof}
Let $f\in R$, $g\in R$. Assume $f\notin I$ and $g\notin I$. Then by
the maximality of $I$, both $I+(f)$ and $I+(g)$ intersect $U$.
There exists elements of the form $af+i\in U$ and $bg+j\in U$ where
$i,j\in I$. If $fg\in I$, then the product $(af+i)(bg+j)\in I$,
contradicting the fact that $I$ is disjoint with $U$.
\end{proof}

\begin{proposition}[Complement of a prime ideal is a multiplicative set]\label{prop:commutative-algebra:complement-of-prime-ideal-is-multiplicative-set}
Let $R$ be a ring. Let $P\in\Spec(R)$ be a prime ideal.
Then $U=R\setminus P$ is a multiplicative set.
\end{proposition}

\begin{proof}
  We see that $1\notin P$. Then
  \begin{equation}\label{eq:complement-of-prime-ideal:condition-one}
1\in U\setminus P.
  \end{equation}
  Let $u\in R\setminus P$ and $v\in R\setminus P$.
  Then $u\notin P$ and $v\notin P$. Then $uv\notin P$, otherwise
  either $u\in P$ or $v\in P$ by definition of a prime ideal. Then
  \begin{equation}
uv\in R\setminus P.
  \end{equation}
  Hence $U$ satisfies the properties of being multiplicatively-closed
  by Equation~\eqref{eq:complement-of-prime-ideal:condition-one} and
  the Definition of multiplicatively-closed sets.
\end{proof}

\begin{definition}[Localization of a ring over a multiplicative set]\mml{ringfrac:def 18}
Let $R$ be any commutative ring.
Let $S$ be a nonempty multiplicatively-closed subset of $R$.
We define the \define{Localization} of $R$ with respect to $S$ to be
the ring $S^{-1}R$ consisting of equivalence classes of fractions
$a/s$ where $a\in R$ and $s\in S$, the equivalence relation being
\begin{equation}
\frac{a}{s}\sim\frac{b}{s'}\iff\exists s''\in S\ldotp s''(s'a-sb)=0,
\end{equation}
where the $s''$ factor is needed (if we omit it, we would not have an
equivalence relation, because of zero-divisors).
\end{definition}

\begin{definition}[Localization with respect to a prime]
Let $R$ be a commutative ring. Let $P\in\Spec(R)$ be a prime ideal.
We define the \define{Localization with respect to the prime} $P$ of
the ring $R$ to be the ring $R_{P}:=R/(R\setminus P)$. We can see this
is well-defined because $U=R\setminus P$ is a multiplicative set by
Proposition~\ref{prop:commutative-algebra:complement-of-prime-ideal-is-multiplicative-set}.
\end{definition}

\section{Nakayama's Lemma}

\begin{remark}[Usefulness of Nakayama's lemma]
You'd use Nakayama's lemma when $R$ is a local ring, and $I\ideal R$
is its maximal ideal.
\end{remark}

\begin{proposition}
Let $R$ be a unital commutative ring. Let $I\ideal R$ be an ideal.
Let $M$ be a finitely-generated module over $R$.
If $IM=M$, then there exists an $r\in R$ with $r\equiv1\bmod{I}$
such that $rM=0$.
\end{proposition}

\begin{proposition}[$IM=M\implies im=m$]
Let $I\ideal R$ be an ideal, let $M$ be a finitely-generated module
over $R$.
If $IM=M$, then there exists an $i\in I$ such that $im=m$ for all
$m\in M$.
\end{proposition}

\begin{proposition}
Let $R$ be a ring, let $M$ be a finitely-generated module over $R$.
Let $I\ideal R$ be an ideal. If $IM=M$, then there exists some $r\in R$
of the form $r=1+x$ where $x\in I$ and $rM=0$. Moreover, if
$I\subset\JacobsonRadical(A)$, then $M=0$.
\end{proposition}

This is how Matsumara~\cite[1M]{matsumura1970commutative} presents
Nakayama's lemma.

\begin{proposition}
Let $M$ is a finitely-generated module over $R$ and
let $\JacobsonRadical(R)$ be the Jacobson radical of $R$.
If $\JacobsonRadical(R)M=M$, then $M=0$.
\end{proposition}